% --- Archon physics formalization source begin ---
% archon:physics
% archon:covers PhyXMiniProblems/problem_phyx_mini_0469.lean
% archon:source-report reports/phyx_mini/problem_phyx_mini_0469.source.json

\chapter{Physics problem phyx\_mini\_0469}
\label{ch:PhyXMiniProblems_problem_phyx_mini_0469}

\paragraph{Problem source.}
\#\# Physical scenario

In an automobile engine, a mixture of air and vaporized gasoline is compressed in the cylinders before being ignited. A typical engine has a compression ratio of 9.00 to 1; that is, the gas in the cylinders is compressed to \textbackslash{}frac\{1\}\{9.00\} of its original volume (figure). The intake and exhaust valves are closed during the compression, so the quantity of gas is constant.Its initial temperature is 27\textasciicircum{}\textbackslash{}circ \textbackslash{}mathrm\{C\} and the initial and final pressures are 1.00 \textbackslash{}, \textbackslash{}text\{atm\} and 21.7 \textbackslash{}, \textbackslash{}text\{atm\}, respectively.

\#\# Question

What is the final temperature of the compressed gas?

\#\# Answer choices

- A: A: 420°C

- B: B: 450°C

- C: C: 530°C

- D: D: 500°C

\#\# Figure caption (auxiliary; use the image as primary evidence)

The image is a cross-sectional diagram of an internal combustion engine. It illustrates various components and their placement within the engine. Key features include:

- **Intake Valve**: Shown at the top of the cylinders, these valves allow air or an air-fuel mixture into the combustion chamber.
- **Exhaust Valve**: Also positioned at the top, these valves release exhaust gases after combustion.
- **Fuel Injector**: Located near the top of each cylinder, responsible for spraying fuel into the combustion chamber.
- **Combustion Chamber**: The enclosed area within the cylinder where fuel combustion occurs.
- **Fuel Pump**: Indicated by a label, this component supplies fuel to the fuel injectors.
- **Pistons**: Illustrated as cylindrical devices within the cylinders, moving up and down to convert fuel energy into mechanical motion.
- **Gear Assembly**: Visible on the left, indicating parts of the timing or drive mechanism.

The labels point to their respective parts, providing an educational overview of engine components and function.

\#\# Dataset metadata

Ideal Gases and Kinetic Theory; Physical Model Grounding Reasoning, Multi-Formula Reasoning

\paragraph{Recorded answer/context.}
B: B: 450°C

\paragraph{Figure/image path.}
/root/proposal\_for\_physic/science-mango/phyx\_mini\_run/phyx\_data/test\_image/469.png

\paragraph{Formalization target.}
create a compiling Lean file with sorry bodies at `PhyXMiniProblems/problem\_phyx\_mini\_0469.lean`. The Lean declarations must preserve the physical quantities, dimensions or dimensional roles, figure labels, governing-law hypotheses, and final relation expressed by this problem.
Use Mathlib/Physlib names found through LeanExplore where available. If a physics API is missing, introduce faithful local abstractions rather than scalar placeholder aliases.

\begin{theorem}[Physics formalization target]
\label{thm:physics:phyx_mini_0469:target}
\lean{PhyXMiniProblems.ProblemPhyXMini0469.problem_phyx_mini_0469}
\uses{def:physics:phyx-mini-0469:phyxminiproblems-problemphyxmini0469-automobileenginecompressionsetup, def:physics:phyx-mini-0469:phyxminiproblems-problemphyxmini0469-statetemperatureinkelvin, def:physics:phyx-mini-0469:phyxminiproblems-problemphyxmini0469-statetemperatureindegreescelsius, def:physics:phyx-mini-0469:phyxminiproblems-problemphyxmini0469-matchesautomobileenginecompressionscenario, def:physics:phyx-mini-0469:phyxminiproblems-problemphyxmini0469-matchesproblemreadouts, def:physics:phyx-mini-0469:phyxminiproblems-problemphyxmini0469-matchessuppliedenginefigure, def:physics:phyx-mini-0469:phyxminiproblems-problemphyxmini0469-hasphysicalcompressionparameters, def:physics:phyx-mini-0469:phyxminiproblems-problemphyxmini0469-satisfiesclosedcylinderamountconservation, def:physics:phyx-mini-0469:phyxminiproblems-problemphyxmini0469-satisfiesendpointidealgaslaw, def:physics:phyx-mini-0469:phyxminiproblems-problemphyxmini0469-roundstodisplayedtendegrees, def:physics:phyx-mini-0469:phyxminiproblems-problemphyxmini0469-isuniqueclosestdisplayedtemperature}
The assigned autoformalize agent should translate this physics problem into Lean declarations in the covered file, with theorem and lemma proof bodies written as `by sorry`.
\end{theorem}
\begin{proof}
This is an autoformalization task, not a proof task. Produce statements that can later be proved by the physics prover without weakening the physical model.
\end{proof}

% --- Archon declaration topology begin ---
\section{Declaration topology}
\begin{definition}[Volume Quantity]
  \label{def:physics:phyx-mini-0469:phyxminiproblems-problemphyxmini0469-volumequantity}
  \lean{PhyXMiniProblems.ProblemPhyXMini0469.VolumeQuantity}
  A nonnegative physical volume carrying the dimension L³
\end{definition}
\begin{proof}
  This is the declared model definition of Volume Quantity.
\end{proof}
\begin{definition}[pressure In Pascals]
  \label{def:physics:phyx-mini-0469:phyxminiproblems-problemphyxmini0469-pressureinpascals}
  \lean{PhyXMiniProblems.ProblemPhyXMini0469.pressureInPascals}
  Read a physical pressure in pascals
\end{definition}
\begin{proof}
  This is the declared model definition of pressure In Pascals.
\end{proof}
\begin{definition}[pressure In Atmospheres]
  \label{def:physics:phyx-mini-0469:phyxminiproblems-problemphyxmini0469-pressureinatmospheres}
  \lean{PhyXMiniProblems.ProblemPhyXMini0469.pressureInAtmospheres}
  \uses{def:physics:phyx-mini-0469:phyxminiproblems-problemphyxmini0469-pressureinpascals}
  Read a physical pressure in standard atmospheres by comparison with Physlib's dimensionful DimPressure.standardAtmosphere
\end{definition}
\begin{proof}
  This is the declared model definition of pressure In Atmospheres.
\end{proof}
\begin{definition}[volume In Cubic Meters]
  \label{def:physics:phyx-mini-0469:phyxminiproblems-problemphyxmini0469-volumeincubicmeters}
  \lean{PhyXMiniProblems.ProblemPhyXMini0469.volumeInCubicMeters}
  \uses{def:physics:phyx-mini-0469:phyxminiproblems-problemphyxmini0469-volumequantity}
  Read a physical volume in cubic metres
\end{definition}
\begin{proof}
  This is the declared model definition of volume In Cubic Meters.
\end{proof}
\begin{definition}[temperature In Kelvin]
  \label{def:physics:phyx-mini-0469:phyxminiproblems-problemphyxmini0469-temperatureinkelvin}
  \lean{PhyXMiniProblems.ProblemPhyXMini0469.temperatureInKelvin}
  Read an absolute temperature in kelvins. Physlib's Temperature stores a nonnegative magnitude in an arbitrary zero-preserving unit, so the storage unit remains explicit in the setup
\end{definition}
\begin{proof}
  This is the declared model definition of temperature In Kelvin.
\end{proof}
\begin{definition}[temperature In Degrees Celsius]
  \label{def:physics:phyx-mini-0469:phyxminiproblems-problemphyxmini0469-temperatureindegreescelsius}
  \lean{PhyXMiniProblems.ProblemPhyXMini0469.temperatureInDegreesCelsius}
  \uses{def:physics:phyx-mini-0469:phyxminiproblems-problemphyxmini0469-temperatureinkelvin}
  The affine Celsius readout attached to an absolute temperature. The offset is 273.15 K = 5463/20 K; Celsius is not treated as an absolute scale
\end{definition}
\begin{proof}
  This is the declared model definition of temperature In Degrees Celsius.
\end{proof}
\begin{definition}[Compression State]
  \label{def:physics:phyx-mini-0469:phyxminiproblems-problemphyxmini0469-compressionstate}
  \lean{PhyXMiniProblems.ProblemPhyXMini0469.CompressionState}
  The two equilibrium states before and after compression
\end{definition}
\begin{proof}
  This is the declared model definition of Compression State.
\end{proof}
\begin{definition}[Gas Model]
  \label{def:physics:phyx-mini-0469:phyxminiproblems-problemphyxmini0469-gasmodel}
  \lean{PhyXMiniProblems.ProblemPhyXMini0469.GasModel}
  Thermodynamic model selected for the cylinder contents
\end{definition}
\begin{proof}
  This is the declared model definition of Gas Model.
\end{proof}
\begin{definition}[Cylinder Contents]
  \label{def:physics:phyx-mini-0469:phyxminiproblems-problemphyxmini0469-cylindercontents}
  \lean{PhyXMiniProblems.ProblemPhyXMini0469.CylinderContents}
  Composition of the gas in the cylinder
\end{definition}
\begin{proof}
  This is the declared model definition of Cylinder Contents.
\end{proof}
\begin{definition}[Valve Position]
  \label{def:physics:phyx-mini-0469:phyxminiproblems-problemphyxmini0469-valveposition}
  \lean{PhyXMiniProblems.ProblemPhyXMini0469.ValvePosition}
  Position of an intake or exhaust valve during the compression stroke
\end{definition}
\begin{proof}
  This is the declared model definition of Valve Position.
\end{proof}
\begin{definition}[Engine Figure Label]
  \label{def:physics:phyx-mini-0469:phyxminiproblems-problemphyxmini0469-enginefigurelabel}
  \lean{PhyXMiniProblems.ProblemPhyXMini0469.EngineFigureLabel}
  Literal component labels printed in the supplied engine raster
\end{definition}
\begin{proof}
  This is the declared model definition of Engine Figure Label.
\end{proof}
\begin{definition}[Engine Figure Object]
  \label{def:physics:phyx-mini-0469:phyxminiproblems-problemphyxmini0469-enginefigureobject}
  \lean{PhyXMiniProblems.ProblemPhyXMini0469.EngineFigureObject}
  Physical engine components visibly depicted in the cross section
\end{definition}
\begin{proof}
  This is the declared model definition of Engine Figure Object.
\end{proof}
\begin{definition}[Engine Cross Section Figure]
  \label{def:physics:phyx-mini-0469:phyxminiproblems-problemphyxmini0469-enginecrosssectionfigure}
  \lean{PhyXMiniProblems.ProblemPhyXMini0469.EngineCrossSectionFigure}
  \uses{def:physics:phyx-mini-0469:phyxminiproblems-problemphyxmini0469-enginefigurelabel, def:physics:phyx-mini-0469:phyxminiproblems-problemphyxmini0469-enginefigureobject}
  Qualitative information retained from the primary raster. The geometric fields describe component placement and connectivity; they contain no pressure, volume, or temperature readout
\end{definition}
\begin{proof}
  This is the declared model definition of Engine Cross Section Figure.
\end{proof}
\begin{definition}[Cylinder Gas State]
  \label{def:physics:phyx-mini-0469:phyxminiproblems-problemphyxmini0469-cylindergasstate}
  \lean{PhyXMiniProblems.ProblemPhyXMini0469.CylinderGasState}
  \uses{def:physics:phyx-mini-0469:phyxminiproblems-problemphyxmini0469-volumequantity}
  Pressure, volume, and absolute temperature of one endpoint gas state
\end{definition}
\begin{proof}
  This is the declared model definition of Cylinder Gas State.
\end{proof}
\begin{definition}[Automobile Engine Compression Setup]
  \label{def:physics:phyx-mini-0469:phyxminiproblems-problemphyxmini0469-automobileenginecompressionsetup}
  \lean{PhyXMiniProblems.ProblemPhyXMini0469.AutomobileEngineCompressionSetup}
  \uses{def:physics:phyx-mini-0469:phyxminiproblems-problemphyxmini0469-compressionstate, def:physics:phyx-mini-0469:phyxminiproblems-problemphyxmini0469-gasmodel, def:physics:phyx-mini-0469:phyxminiproblems-problemphyxmini0469-cylindercontents, def:physics:phyx-mini-0469:phyxminiproblems-problemphyxmini0469-valveposition, def:physics:phyx-mini-0469:phyxminiproblems-problemphyxmini0469-enginecrosssectionfigure, def:physics:phyx-mini-0469:phyxminiproblems-problemphyxmini0469-cylindergasstate}
  The physical setup. GasAmount remains abstract and has a separate mole readout, so amount of substance is not identified with a bare scalar type. The state-dependent amount makes conservation during the closed-valve stroke an explicit law rather than a fact hidden by construction
\end{definition}
\begin{proof}
  This is the declared model definition of Automobile Engine Compression Setup.
\end{proof}
\begin{definition}[state Pressure In Atmospheres]
  \label{def:physics:phyx-mini-0469:phyxminiproblems-problemphyxmini0469-statepressureinatmospheres}
  \lean{PhyXMiniProblems.ProblemPhyXMini0469.statePressureInAtmospheres}
  \uses{def:physics:phyx-mini-0469:phyxminiproblems-problemphyxmini0469-pressureinatmospheres, def:physics:phyx-mini-0469:phyxminiproblems-problemphyxmini0469-compressionstate, def:physics:phyx-mini-0469:phyxminiproblems-problemphyxmini0469-automobileenginecompressionsetup}
  Pressure in atmospheres at a named endpoint
\end{definition}
\begin{proof}
  This is the declared model definition of state Pressure In Atmospheres.
\end{proof}
\begin{definition}[state Volume In Cubic Meters]
  \label{def:physics:phyx-mini-0469:phyxminiproblems-problemphyxmini0469-statevolumeincubicmeters}
  \lean{PhyXMiniProblems.ProblemPhyXMini0469.stateVolumeInCubicMeters}
  \uses{def:physics:phyx-mini-0469:phyxminiproblems-problemphyxmini0469-volumeincubicmeters, def:physics:phyx-mini-0469:phyxminiproblems-problemphyxmini0469-compressionstate, def:physics:phyx-mini-0469:phyxminiproblems-problemphyxmini0469-automobileenginecompressionsetup}
  Volume in cubic metres at a named endpoint
\end{definition}
\begin{proof}
  This is the declared model definition of state Volume In Cubic Meters.
\end{proof}
\begin{definition}[state Temperature In Kelvin]
  \label{def:physics:phyx-mini-0469:phyxminiproblems-problemphyxmini0469-statetemperatureinkelvin}
  \lean{PhyXMiniProblems.ProblemPhyXMini0469.stateTemperatureInKelvin}
  \uses{def:physics:phyx-mini-0469:phyxminiproblems-problemphyxmini0469-temperatureinkelvin, def:physics:phyx-mini-0469:phyxminiproblems-problemphyxmini0469-compressionstate, def:physics:phyx-mini-0469:phyxminiproblems-problemphyxmini0469-automobileenginecompressionsetup}
  Absolute temperature in kelvins at a named endpoint
\end{definition}
\begin{proof}
  This is the declared model definition of state Temperature In Kelvin.
\end{proof}
\begin{definition}[state Temperature In Degrees Celsius]
  \label{def:physics:phyx-mini-0469:phyxminiproblems-problemphyxmini0469-statetemperatureindegreescelsius}
  \lean{PhyXMiniProblems.ProblemPhyXMini0469.stateTemperatureInDegreesCelsius}
  \uses{def:physics:phyx-mini-0469:phyxminiproblems-problemphyxmini0469-temperatureindegreescelsius, def:physics:phyx-mini-0469:phyxminiproblems-problemphyxmini0469-compressionstate, def:physics:phyx-mini-0469:phyxminiproblems-problemphyxmini0469-automobileenginecompressionsetup}
  Celsius temperature at a named endpoint
\end{definition}
\begin{proof}
  This is the declared model definition of state Temperature In Degrees Celsius.
\end{proof}
\begin{definition}[Matches Automobile Engine Compression Scenario]
  \label{def:physics:phyx-mini-0469:phyxminiproblems-problemphyxmini0469-matchesautomobileenginecompressionscenario}
  \lean{PhyXMiniProblems.ProblemPhyXMini0469.MatchesAutomobileEngineCompressionScenario}
  \uses{def:physics:phyx-mini-0469:phyxminiproblems-problemphyxmini0469-automobileenginecompressionsetup}
  Qualitative physical information from the prose. Closed valves are recorded here, while their consequence for gas amount is stated separately as a process law
\end{definition}
\begin{proof}
  This is the declared model definition of Matches Automobile Engine Compression Scenario.
\end{proof}
\begin{definition}[Matches Problem Readouts]
  \label{def:physics:phyx-mini-0469:phyxminiproblems-problemphyxmini0469-matchesproblemreadouts}
  \lean{PhyXMiniProblems.ProblemPhyXMini0469.MatchesProblemReadouts}
  \uses{def:physics:phyx-mini-0469:phyxminiproblems-problemphyxmini0469-automobileenginecompressionsetup, def:physics:phyx-mini-0469:phyxminiproblems-problemphyxmini0469-statepressureinatmospheres, def:physics:phyx-mini-0469:phyxminiproblems-problemphyxmini0469-statevolumeincubicmeters, def:physics:phyx-mini-0469:phyxminiproblems-problemphyxmini0469-statetemperatureindegreescelsius}
  The four numerical readouts printed in the problem. The final temperature and all answer choices are absent. 9.00 : 1 is expressed as V\_final = V\_initial / 9
\end{definition}
\begin{proof}
  This is the declared model definition of Matches Problem Readouts.
\end{proof}
\begin{definition}[Matches Supplied Engine Figure]
  \label{def:physics:phyx-mini-0469:phyxminiproblems-problemphyxmini0469-matchessuppliedenginefigure}
  \lean{PhyXMiniProblems.ProblemPhyXMini0469.MatchesSuppliedEngineFigure}
  \uses{def:physics:phyx-mini-0469:phyxminiproblems-problemphyxmini0469-enginecrosssectionfigure}
  Facts read directly from the supplied cross-sectional engine raster
\end{definition}
\begin{proof}
  This is the declared model definition of Matches Supplied Engine Figure.
\end{proof}
\begin{definition}[Has Physical Compression Parameters]
  \label{def:physics:phyx-mini-0469:phyxminiproblems-problemphyxmini0469-hasphysicalcompressionparameters}
  \lean{PhyXMiniProblems.ProblemPhyXMini0469.HasPhysicalCompressionParameters}
  \uses{def:physics:phyx-mini-0469:phyxminiproblems-problemphyxmini0469-automobileenginecompressionsetup, def:physics:phyx-mini-0469:phyxminiproblems-problemphyxmini0469-statepressureinatmospheres, def:physics:phyx-mini-0469:phyxminiproblems-problemphyxmini0469-statevolumeincubicmeters, def:physics:phyx-mini-0469:phyxminiproblems-problemphyxmini0469-statetemperatureinkelvin}
  Positivity and nondegeneracy conditions for both gas states
\end{definition}
\begin{proof}
  This is the declared model definition of Has Physical Compression Parameters.
\end{proof}
\begin{definition}[Satisfies Closed Cylinder Amount Conservation]
  \label{def:physics:phyx-mini-0469:phyxminiproblems-problemphyxmini0469-satisfiesclosedcylinderamountconservation}
  \lean{PhyXMiniProblems.ProblemPhyXMini0469.SatisfiesClosedCylinderAmountConservation}
  \uses{def:physics:phyx-mini-0469:phyxminiproblems-problemphyxmini0469-automobileenginecompressionsetup}
  Closed intake and exhaust valves make the amount of gas at the two endpoint states equal. This is a process law, not a final-temperature relation
\end{definition}
\begin{proof}
  This is the declared model definition of Satisfies Closed Cylinder Amount Conservation.
\end{proof}
\begin{definition}[Satisfies Endpoint Ideal Gas Law]
  \label{def:physics:phyx-mini-0469:phyxminiproblems-problemphyxmini0469-satisfiesendpointidealgaslaw}
  \lean{PhyXMiniProblems.ProblemPhyXMini0469.SatisfiesEndpointIdealGasLaw}
  \uses{def:physics:phyx-mini-0469:phyxminiproblems-problemphyxmini0469-automobileenginecompressionsetup, def:physics:phyx-mini-0469:phyxminiproblems-problemphyxmini0469-statepressureinatmospheres, def:physics:phyx-mini-0469:phyxminiproblems-problemphyxmini0469-statevolumeincubicmeters, def:physics:phyx-mini-0469:phyxminiproblems-problemphyxmini0469-statetemperatureinkelvin}
  The ideal-gas equation pV = nRT, applied uniformly at both endpoints in coherent atmosphere, cubic-metre, mole, and kelvin readouts. It does not assign the unknown final temperature or select an answer choice. Physlib's IdealGas.ideal\_gas\_law is specialized to its unitless statistical mechanics model with R = 1; this local law preserves the dimensional engine state and explicit unit readouts needed here
\end{definition}
\begin{proof}
  This is the declared model definition of Satisfies Endpoint Ideal Gas Law.
\end{proof}
\begin{definition}[Answer Choice]
  \label{def:physics:phyx-mini-0469:phyxminiproblems-problemphyxmini0469-answerchoice}
  \lean{PhyXMiniProblems.ProblemPhyXMini0469.AnswerChoice}
  Labels of the four choices printed with the problem
\end{definition}
\begin{proof}
  This is the declared model definition of Answer Choice.
\end{proof}
\begin{definition}[displayed Temperature In Degrees Celsius]
  \label{def:physics:phyx-mini-0469:phyxminiproblems-problemphyxmini0469-displayedtemperatureindegreescelsius}
  \lean{PhyXMiniProblems.ProblemPhyXMini0469.displayedTemperatureInDegreesCelsius}
  \uses{def:physics:phyx-mini-0469:phyxminiproblems-problemphyxmini0469-answerchoice}
  Celsius value displayed beside each answer label
\end{definition}
\begin{proof}
  This is the declared model definition of displayed Temperature In Degrees Celsius.
\end{proof}
\begin{definition}[recorded Dataset Answer]
  \label{def:physics:phyx-mini-0469:phyxminiproblems-problemphyxmini0469-recordeddatasetanswer}
  \lean{PhyXMiniProblems.ProblemPhyXMini0469.recordedDatasetAnswer}
  \uses{def:physics:phyx-mini-0469:phyxminiproblems-problemphyxmini0469-answerchoice}
  Dataset metadata recording the supplied answer label
\end{definition}
\begin{proof}
  This is the declared model definition of recorded Dataset Answer.
\end{proof}
\begin{definition}[Rounds To Displayed Ten Degrees]
  \label{def:physics:phyx-mini-0469:phyxminiproblems-problemphyxmini0469-roundstodisplayedtendegrees}
  \lean{PhyXMiniProblems.ProblemPhyXMini0469.RoundsToDisplayedTenDegrees}
  \uses{def:physics:phyx-mini-0469:phyxminiproblems-problemphyxmini0469-automobileenginecompressionsetup, def:physics:phyx-mini-0469:phyxminiproblems-problemphyxmini0469-statetemperatureindegreescelsius, def:physics:phyx-mini-0469:phyxminiproblems-problemphyxmini0469-answerchoice, def:physics:phyx-mini-0469:phyxminiproblems-problemphyxmini0469-displayedtemperatureindegreescelsius}
  A computed temperature matches a choice printed to the nearest ten degrees when it lies within five degrees of that displayed value
\end{definition}
\begin{proof}
  This is the declared model definition of Rounds To Displayed Ten Degrees.
\end{proof}
\begin{definition}[Is Unique Closest Displayed Temperature]
  \label{def:physics:phyx-mini-0469:phyxminiproblems-problemphyxmini0469-isuniqueclosestdisplayedtemperature}
  \lean{PhyXMiniProblems.ProblemPhyXMini0469.IsUniqueClosestDisplayedTemperature}
  \uses{def:physics:phyx-mini-0469:phyxminiproblems-problemphyxmini0469-automobileenginecompressionsetup, def:physics:phyx-mini-0469:phyxminiproblems-problemphyxmini0469-statetemperatureindegreescelsius, def:physics:phyx-mini-0469:phyxminiproblems-problemphyxmini0469-answerchoice, def:physics:phyx-mini-0469:phyxminiproblems-problemphyxmini0469-displayedtemperatureindegreescelsius}
  The selected displayed temperature is strictly closer than every rival
\end{definition}
\begin{proof}
  This is the declared model definition of Is Unique Closest Displayed Temperature.
\end{proof}
\begin{lemma}[final temperature exact]
  \label{lem:physics:phyx-mini-0469:phyxminiproblems-problemphyxmini0469-final-temperature-exact}
  \lean{PhyXMiniProblems.ProblemPhyXMini0469.final_temperature_exact}
  \uses{def:physics:phyx-mini-0469:phyxminiproblems-problemphyxmini0469-automobileenginecompressionsetup, def:physics:phyx-mini-0469:phyxminiproblems-problemphyxmini0469-statetemperatureinkelvin, def:physics:phyx-mini-0469:phyxminiproblems-problemphyxmini0469-statetemperatureindegreescelsius, def:physics:phyx-mini-0469:phyxminiproblems-problemphyxmini0469-matchesproblemreadouts, def:physics:phyx-mini-0469:phyxminiproblems-problemphyxmini0469-hasphysicalcompressionparameters, def:physics:phyx-mini-0469:phyxminiproblems-problemphyxmini0469-satisfiesclosedcylinderamountconservation, def:physics:phyx-mini-0469:phyxminiproblems-problemphyxmini0469-satisfiesendpointidealgaslaw}
  The endpoint ideal-gas equations, conservation of gas amount, and supplied ratios give the exact model value T\_final = 144739/200 K = 90109/200 °C = 450.545 °C
\end{lemma}
\begin{proof}
  The conclusion follows from the governing relations and figure readouts named in the statement of final temperature exact.
\end{proof}
% --- Archon declaration topology end ---
% --- Archon physics formalization source end ---
